\documentclass{article}

% Use numeric citations (required for thebibliography without bibtex)
\PassOptionsToPackage{numbers,compress}{natbib}

% Workshop: single-blind reviewing
\usepackage[sglblindworkshop]{neurips_2026}
\workshoptitle{Foundations of Agentic Systems Theory (FAST @ NeurIPS 2026)}

\usepackage[utf8]{inputenc}
\usepackage[T1]{fontenc}
\usepackage{hyperref}
\usepackage{url}
\usepackage{booktabs}
\usepackage{amsfonts}
\usepackage{amsmath}
\usepackage{amssymb}
\usepackage{amsthm}
\usepackage{nicefrac}
\usepackage{microtype}
\usepackage{xcolor}
\usepackage{graphicx}
\usepackage{multirow}
\usepackage{listings}

\newtheorem{definition}{Definition}
\newtheorem{proposition}{Proposition}
\newtheorem{theorem}{Theorem}
\newtheorem{corollary}{Corollary}

\lstset{
  basicstyle=\ttfamily\small,
  breaklines=true,
  frame=single,
  xleftmargin=1em,
}

\title{Fail-Closed Agentic Orchestration: A Lattice-Theoretic Characterization of Safe State Reducers, with a Commerce Instantiation}

\author{%
  Anonymous Author(s)
}

\begin{document}

\maketitle

\begin{abstract}
LLM orchestration frameworks---LangGraph, AutoGen, OpenAI Agents SDK---give no formal safety
guarantee for consequential financial actions: any node can overwrite a prior fraud-block
decision. We show this is structural, and characterize the fix. For any bounded safety lattice
$(L,\le,\bot,\top)$, call a channel reducer \emph{safe} if it is fail-closed, non-regressive,
order-independent, and replay-stable. Our main result (Theorem~\ref{thm:general_char}) is that
a reducer is safe \emph{iff} it is a semilattice operation dominating the lattice join. Hence
the join is the unique pointwise-minimal safe reducer---any other safe reducer escalates
strictly more often without adding safety---and last-writer-wins is unsafe on every lattice
with at least two levels (Corollary~\ref{cor:lww}). A delegation-closure theorem
(Theorem~\ref{thm:delegation}) lifts this to hierarchies: sub-agent decisions compose safely
into a parent channel exactly when the embedding between their lattices is top-preserving, so
an agent population is fail-closed as a whole.

We instantiate the theory as \textbf{Commerce-Native Agentic Orchestration (CAO)}, providing
five transaction properties---Execution, Atomicity, Effectiveness, Verification, Traceability
(EAVT)---via Commerce-Typed State on $\{\textit{allow}<\textit{flag}<\textit{block}\}$ (where
the minimal safe reducer is \texttt{block\_wins}), an ISO~8583-grounded automaton with $O(n)$
decidable validity, ceremony tiers, a structural atomic gate, and a hash-chained audit log.
Empirically the verification cascade reaches 88\% behavioral recall at 0\% hard false-positive
rate over ten attack classes at $<$\$0.01/1k sessions, versus 61\% for a full-session Claude
Sonnet~4.6 judge at \$3/1k. Code and results are open-source; benchmark details
in~\cite{acpbench}.
\end{abstract}

\section{Introduction}

LLM agents are becoming \emph{consequential actors}: systems authorized to book travel,
subscribe to software, and pay for API services without per-action human approval. This
demands a formal account of what it means for an agent to take a financial action, not merely
to route between graph nodes. Distributed systems and payment engineering have supplied such
accounts for four decades---Saga~\cite{saga} for long-running atomicity, two-phase
commit~\cite{2pc} for consensus before irreversible commitment, ISO~8583~\cite{iso8583} for
the financial message lifecycle, Certificate Transparency~\cite{ct} for tamper-evident logs.
These are load-bearing invariants for commerce systems, and no existing LLM orchestration
framework enforces any of them.

\paragraph{Five transaction properties.} Sound commerce orchestration must provide
\textbf{Execution} (the agent follows a valid state grammar $\mathcal{M}$; violations such as
START$\to$COMMIT are caught before the rail executes), \textbf{Atomicity} (actions commit
all-or-nothing via risk-adjusted ceremony tiers; a RESERVE-then-failure produces a
compensating VOID), \textbf{Effectiveness} (no path reaches the payment rail bypassing the
gate), \textbf{Verification} (a confidence cascade gives risk-adjusted verification depth),
and \textbf{Traceability} (a tamper-evident audit trail per action)---\emph{EAVT}.

\paragraph{The framework gap.} Table~\ref{tab:framework_compare} evaluates existing
orchestration frameworks against these five properties; all were designed for graph routing and
multi-agent coordination, not for enforcing the invariants financial commitment requires. The
``partial'' entries mark mechanisms that exist syntactically but do not enforce the safety
semantic: LangGraph's \texttt{interrupt\_before} gates workflow without any protocol-validity
check (and LWW still lets a later node overwrite a block, Thm.~\ref{thm:lww_unsound}), while
LangSmith logs observability data without hash-chain tamper-evidence or third-party-verifiable
inclusion proofs; OpenAI Agents SDK guardrails observe text only, enforcing no action-sequence
grammar and producing no audit log.

\begin{table}[h]
\caption{Five transaction properties across orchestration frameworks. CAO provides all five at the semantic level.}
\label{tab:framework_compare}
\centering
\footnotesize
\begin{tabular}{p{3.2cm}cccc}
\toprule
Property & LangGraph & AutoGen & OpenAI Agents & CAO \\
\midrule
Execution (grammar $\mathcal{M}$) & $\times$ & $\times$ & $\times$ & \checkmark \\
Atomicity (Saga/2PC) & $\times$ & $\times$ & $\times$ & \checkmark \\
Effectiveness (gate) & partial & $\times$ & partial & \checkmark \\
Verification (cascade) & $\times$ & $\times$ & partial & \checkmark \\
Traceability (CARLog) & partial & $\times$ & $\times$ & \checkmark \\
\bottomrule
\end{tabular}
\end{table}

\paragraph{The three-layer stack.} Agentic commerce safety is defense-in-depth across three
layers. \emph{Layer~0} (agent reasoning: LLM, tools, memory, planning) is where garak,
promptfoo, and LlamaGuard observe text output and tool selection; it is the right place to
catch a jailbroken agent, but it cannot see the protocol wire or wallet state.
\emph{Layer~1} (the MCP/protocol wire, this paper) carries typed
\texttt{ConsequentialAction} fields---action type, merchant, MCC, amount, idempotency key,
wallet state---invisible at Layer~0, and is the right place to catch a
\emph{correctly-reasoning} agent whose action violates a payment invariant (COMMIT before
QUOTE, replayed idempotency key, out-of-allowlist MCC, ceiling breach). \emph{Layer~2+}
(Stripe Radar, issuer authorization) fires only after the rail call, too late for
in-session recovery or a compensating void. CAO formalizes Layer~1.

\paragraph{Contributions.} None of CAO's individual mechanisms is new---DFA simulation,
Saga/2PC, CRDT joins, and hash-chain logs each have decades of literature. The contributions are:

\begin{enumerate}
\item \textbf{A characterization of safe state-channel reducers}
(\S\ref{sec:general_typing}, Theorem~\ref{thm:general_char}): for any bounded safety
lattice, a reducer is safe \emph{iff} it is a semilattice operation dominating the join.
The join is therefore the unique pointwise-minimal safe reducer, so it is also
precision-optimal (Corollary~\ref{cor:precision}); last-writer-wins is unsafe on every
non-trivial lattice (Corollary~\ref{cor:lww}). Fail-closed safety is thus a property of
the reducer \emph{algebra}, not of application code.

\item \textbf{Delegation closure for agent hierarchies}
(Theorem~\ref{thm:delegation}): sub-agent decisions compose safely into a parent channel
\emph{iff} the embedding between their lattices is top-preserving---the necessary and
sufficient condition for a population of agents to be fail-closed as a whole.

\item \textbf{The EAVT property set and its CAO instantiation}
(\S\ref{sec:automaton}--\ref{sec:cascade}): Commerce-Typed State on
$\{\textit{allow}<\textit{flag}<\textit{block}\}$, where \texttt{block\_wins} is the minimal
safe reducer (Theorem~\ref{thm:type_safety}) and last-writer-wins provably is not
(Theorem~\ref{thm:lww_unsound}); an ISO~8583-grounded automaton with $O(n)$ decidable validity
(Proposition~\ref{prop:decidability}), ceremony tiers, a structural atomic gate, and a
hash-chained CARLog---together providing all five properties that existing frameworks
provide none of at the semantic level.

\item \textbf{Empirical validation} (\S\ref{sec:experiments}): formal properties verified
by construction and by fault injection; the verification cascade measured on
ACP-Bench~\cite{acpbench} at 88\% behavioral recall and 0\% hard FPR.
\end{enumerate}

\section{Background and Related Work}

\paragraph{Orchestration frameworks.} ReAct~\cite{react} interleaves reasoning and tool calls
with no state management beyond the context window. LangGraph~\cite{langgraph} provides
stateful graph execution over Pregel~\cite{pregel}: state is a typed dictionary with
per-channel reducers, and \texttt{interrupt\_before} pauses for human review. Its reducers are
user-specified and unconstrained---the framework imposes no semantics on them---and
\texttt{interrupt\_before} is a workflow gate, not a fraud gate. AutoGen~\cite{autogen} uses
role-based prompting among conversational actors and does not model a financial action
lifecycle. The OpenAI Agents SDK~\cite{openai_agents} applies guardrails to input and output
text, enforcing no grammar over the action sequence and producing no audit log. All operate at
the reasoning and routing layer; none provides the EAVT properties for financial commitments.

\paragraph{Payment protocol foundations.} ISO~8583~\cite{iso8583} orders message types
0100~$\to$~0200~$\to$~0220~$\to$~0420---exactly the invariant cold-start attacks violate.
x402~\cite{x402} maps this lifecycle onto HTTP headers for AI-to-AI payments and documents the
revert-grant attack. Stripe, Adyen, and ISO~20022 all require idempotency keys, which are
unenforceable if the verifier cannot see the settled-key registry. In every prior system the
state machine was enforced by terminal hardware and issuer logic; an LLM agent calling a rail
API directly can skip states. \textbf{There is no terminal---the agent is the terminal.}

\paragraph{Transaction theory and typed state.} Two-phase commit~\cite{2pc} guarantees no
participant commits if any aborts; Saga~\cite{saga} trades atomicity for availability via
compensating transactions $C_i$. CRDTs~\cite{crdts} are join-semilattice structures whose
operation \emph{encodes} an invariant---a G-Counter's join is \texttt{max}, so it can only
increase. Theorem~\ref{thm:general_char} supplies the converse for safety channels: the CRDT
algebra is not merely sufficient but \emph{forced} by natural operational safety requirements.
Typestate~\cite{typestate} and session types~\cite{sessiontypes} constrain which
\emph{operations} a value or channel may undergo next; CTS is the complementary discipline on
which \emph{values} it may take next, restricting evolution to a monotone path in a safety
lattice. Both encode a behavioral invariant in the type rather than in application logic.

\section{Commerce-Native Agentic Orchestration (CAO)}
\label{sec:cao}

\subsection{General Safety Typing for State Channels}
\label{sec:general_typing}

We first state the framework-independent result. A graph-execution framework
(Pregel~\cite{pregel}, LangGraph~\cite{langgraph}) exposes \emph{state channels}:
a value domain together with a reducer that folds node writes into the channel.
We ask which reducers can carry a safety decision at all.

\begin{definition}[Bounded safety lattice]
A \emph{bounded safety lattice} is a finite lattice $(L,\leq)$ with least element
$\bot$ (``no concern'') and greatest element $\top$ (``blocked''). We write
$\sqcup$ for its join.
\end{definition}

\begin{definition}[Channel reducer]
A \emph{channel reducer} on $L$ is a map $r: L\times L\to L$ with $r(\bot,x)=x$
(an empty channel adopts its first write). Given node writes $b_1,\dots,b_n\in L$,
the channel state is the fold $D_0=\bot$, $D_i=r(D_{i-1},b_i)$.
\end{definition}

\begin{definition}[Safe reducer]
\label{def:safe}
$r$ is \emph{safe} for $L$ if for every finite trace $b_1,\dots,b_n$:
\textbf{(S1)}~\emph{fail-closed}---$b_i=\top\Rightarrow D_j=\top$ for all $j\ge i$;
\textbf{(S2)}~\emph{non-regressive}---$D_{i-1}\le D_i$ and $b_i\le D_i$ for all $i$;
\textbf{(S3)}~\emph{order-independent}---every permutation of the trace yields the same $D_n$,
so concurrent writes within a superstep fold deterministically; and
\textbf{(S4)}~\emph{replay-stable}---duplicating any $b_i$ leaves $D_n$ unchanged, so retries
are idempotent.
\end{definition}

\begin{theorem}[Characterization of safe reducers]
\label{thm:general_char}
Let $L$ be a bounded safety lattice. A channel reducer $r$ is safe for $L$ if and
only if $r$ is a semilattice operation (associative, commutative, idempotent)
that dominates the join: $x\sqcup y\le r(x,y)$ for all $x,y\in L$.
Consequently $\sqcup$ is the unique \emph{pointwise-minimal} safe reducer:
every other safe reducer satisfies $r(x,y)>x\sqcup y$ for some pair, i.e.\ it
escalates strictly more often than $\sqcup$ without adding safety.
\end{theorem}

\begin{proof}[Proof sketch; full proof in Appendix~\ref{app:general_char}]
($\Leftarrow$) A semilattice operation gives (S3) and (S4); domination gives (S2); and (S1)
follows from (S2) since $\top$ is maximal. ($\Rightarrow$) Because $r(\bot,x)=x$, every $x$ is
a reachable state, so the trace conditions become identities on $L$: (S3) on two- and
three-element traces yields commutativity and associativity, (S4) on $(x,x)$ yields
idempotence, and (S2) makes $r(x,y)$ an upper bound of $\{x,y\}$, hence $\ge x\sqcup y$.
Minimality follows since $\sqcup$ is itself safe and pointwise $\le$ every safe $r$. \qed
\end{proof}

\begin{corollary}[Chains; \texttt{block\_wins}]
\label{cor:chain}
If $L$ is a chain then $\sqcup=\max$ is the unique minimal safe reducer. For
$L=\{\textit{allow}<\textit{flag}<\textit{block}\}$, $\max$ is exactly \texttt{block\_wins}.
\end{corollary}

\begin{corollary}[Last-writer-wins is unsafe on every non-trivial lattice]
\label{cor:lww}
$r(x,y)=y$ violates (S2) whenever $|L|\ge2$: $r(\top,\bot)=\bot\not\ge\top$.
No framework whose default channel semantics is last-writer-wins can carry a
fail-closed decision without replacing the reducer.
\end{corollary}

\begin{corollary}[Precision optimality]
\label{cor:precision}
Among safe reducers, $\sqcup$ has an empty spurious-escalation set
$\{(x,y): r(x,y)>x\sqcup y\}$. Any stricter safe reducer only adds escalations
(soft-flag false positives), never safety. This is the formal reason the empirical
escalation rate in \S\ref{sec:experiments} is attributable to the verifiers, not the reducer.
\end{corollary}

\begin{theorem}[Delegation closure]
\label{thm:delegation}
Let a parent graph $P$ carry a safe reducer $r_P$ on $L_P$ and delegate to a
sub-agent $S$ carrying a safe reducer $r_S$ on $L_S$; $S$'s final state $d^S$ is
written back into $P$'s channel through an embedding $\iota:L_S\to L_P$.
The composite is fail-closed---(i) a block in $P$ before delegation persists
regardless of $d^S$, and (ii) a block in $S$ forces $P$ to $\top_P$---if and only
if $\iota(\top_S)=\top_P$. If moreover $\iota$ is monotone, severity order is
preserved across the boundary.
\end{theorem}

\begin{proof}[Proof in Appendix~\ref{app:general_char}]
Both directions follow from (S1) and (S2) applied across the delegation boundary; the
converse exhibits a trace in which a non-top-preserving $\iota$ loses the sub-agent's block. \qed
\end{proof}

Theorem~\ref{thm:delegation} is the compositional-safety statement: an agent
population in which every member uses a safe reducer and every delegation edge is
top-preserving is fail-closed as a whole, regardless of depth or fan-out. The
remainder of this section instantiates the framework for agentic commerce with
$L=\{\textit{allow}<\textit{flag}<\textit{block}\}$.

\subsection{Execution Property: Commerce State Automaton}
\label{sec:automaton}

\begin{definition}[Commerce State Automaton]
$\mathcal{M} = (Q, \Sigma, \delta, q_0, F)$ with
$Q = \{\texttt{START}, \texttt{FIND}, \texttt{QUOTE}, \texttt{RESERVE}, \texttt{COMMIT},$
$\texttt{SETTLE}, \texttt{VOID}, \texttt{END}\}$,
$\Sigma = \{\textit{find},\textit{quote},\textit{reserve},\textit{commit},\textit{settle},\textit{void},\textit{refund}\}$,
$q_0 = \texttt{START}$, $F = \{\texttt{END}\}$, and $\delta$ given by
\texttt{START}$\xrightarrow{\textit{find}}$\texttt{FIND}$\xrightarrow{\textit{find}}$\texttt{FIND}$\xrightarrow{\textit{quote}}$\texttt{QUOTE};
\texttt{QUOTE}$\xrightarrow{\textit{reserve}}$\texttt{RESERVE} (2PC prepare) or
$\xrightarrow{\textit{commit}}$\texttt{COMMIT} (Saga direct);
\texttt{RESERVE}$\xrightarrow{\textit{commit}}$\texttt{COMMIT} or $\xrightarrow{\textit{void}}$\texttt{VOID};
\texttt{COMMIT}$\xrightarrow{\textit{settle}}$\texttt{SETTLE} (ISO~8583 0220) or
$\xrightarrow{\textit{refund}}$\texttt{VOID}; any state $\to$ \texttt{END}.
\end{definition}

\begin{definition}[Valid Session]
\label{def:valid_session}
A trajectory $s = a_1,\ldots,a_n$ is \emph{protocol-valid} iff (1)~its state sequence is
accepted by $\mathcal{M}$; and for every \textit{commit} action $a_i$:
(2)~$a_i.\textit{amt} \le \textit{wallet.per\_transaction}$;
(3)~$a_i.\textit{mcc} \in \textit{wallet.allowed\_mcc}$;
(4)~$a_i$'s idempotency key is absent from the settled-key registry; and
(5)~$\textit{ctx.agent\_id}$ matches the agent registered for $\textit{ctx.session\_id}$.
Every attack class in the companion taxonomy violates exactly one condition
(B1$\to$1, B8$\to$2, B7$\to$3, B5$\to$4, B4$\to$5).
\end{definition}

\begin{proposition}[Decidability and complexity]
\label{prop:decidability}
$L(\mathcal{M})$ is regular, so state-transition validity is decidable by single-pass DFA
simulation in $O(n)$ time and $O(|Q|)$ space. Conditions~2--5 of
Definition~\ref{def:valid_session} are decidable in $O(n)$ given wallet state and the
settled-key and agent-ID registries (constant-time threshold comparisons; $O(1)$ amortized
hash-set membership). Hence full protocol validity is decidable in $O(n)$, and the
protocol-violation problem is in $\mathbf{P}$. Moreover an invalid $s$ admits an $O(n)$-checkable
certificate---the first violating action---supporting lightweight chargeback dispute resolution.
\end{proposition}

\subsection{Verification and Effectiveness: Commerce-Typed State}
\label{sec:cts}

LangGraph state channels are (type, reducer) pairs whose default reducer is
last-writer-wins. Commerce-Typed State (CTS) replaces the reducer on each channel with the
join of an appropriate safety lattice, instantiating \S\ref{sec:general_typing}:

\begin{lstlisting}
spend:      Annotated[Decimal, operator.add]       # monotone: spend only grows
decision:   Annotated[str, block_wins]             # block > flag > allow
provenance: Annotated[list[CARRef], operator.add]  # append-only audit trail
flags:      Annotated[frozenset, frozenset.union]  # flags never removed
\end{lstlisting}

\noindent where \texttt{block\_wins} is $\max$ on $D=\{\texttt{allow}<\texttt{flag}<\texttt{block}\}$.
Each channel's reducer is the join of a bounded lattice ($D$; non-negative reals; the
powerset lattice), so Theorem~\ref{thm:general_char} applies directly:

\begin{theorem}[CTS type safety and minimality]
\label{thm:type_safety}
\label{thm:char}
In a CTS-typed graph, if any verifier writes \texttt{block} at step $i$, then
$\mathrm{decision}(s_j)=\texttt{block}$ for all $j\ge i$, for every execution trace and every
interleaving. Moreover \texttt{block\_wins} is the unique pointwise-minimal safe reducer on
$D$: CTS is the \emph{minimal} typed extension of LangGraph's channel model that is
fail-closed, and every alternative safe reducer strictly increases the escalation rate.
\end{theorem}

\begin{proof}
$D$ is a chain with top \texttt{block}, and \texttt{block\_wins} $=\max=\sqcup_D$. Apply
Theorem~\ref{thm:general_char} (safety) with Corollaries~\ref{cor:chain} and
\ref{cor:precision} (minimality and precision-optimality). \qed
\end{proof}

\begin{theorem}[LWW unsoundness]
\label{thm:lww_unsound}
Let $G$ use last-writer-wins on the decision channel. Then for \emph{every} fraud gate node
$n_1$ that can output \texttt{block}, there is a two-node trace
$n_1\to n_2\to\mathrm{payment\_rail}$ in which $n_1$ outputs \texttt{block} and the rail executes.
\end{theorem}

\begin{proof}
Take $n_2$ to be any node writing $\texttt{decision}=\texttt{allow}$---a reasoning node that
re-evaluates the action positively, or one replaying a cached verdict. Under LWW,
$\mathrm{decision}(s_2)=r(\texttt{block},\texttt{allow})=\texttt{allow}$, and the routing guard
$\mathrm{decision}\ne\texttt{block}$ admits the rail. This is an instance of
Corollary~\ref{cor:lww}: LWW violates (S2), so no gate placed in such a graph can be
fail-closed. The counterexample in \S\ref{sec:gate} realizes this in 15 lines of LangGraph. \qed
\end{proof}

\begin{proposition}[Rules-certainty veto preserves safety]
\label{prop:veto}
The veto (block$\to$flag downgrade when LocalRules score$=0$ and ML uncertainty $>0.20$)
fires at the \emph{action} level, modifying the ACP node's output $b_i$ before the
session-level join. Since $D_i=\texttt{block\_wins}(D_{i-1},b_i)$, if $D_{i-1}=\texttt{block}$
then $D_i=\texttt{block}$ regardless of the veto: it cannot reduce a decision already in the
log (Theorem~\ref{thm:type_safety}). It only converts an \emph{uncertain} ML block into a
reviewable flag when no rule has fired, and is thus a precision mechanism, not a safety bypass.
\end{proposition}

\subsection{Effectiveness: Structural Atomic Gate}
\label{sec:gate}

Theorem~\ref{thm:lww_unsound} is realized concretely below: \texttt{reason\_node} silently
overwrites the gate's block under LWW, and \texttt{pay\_node} executes.

\begin{lstlisting}
class State(TypedDict):
    decision: str          # last-writer-wins: no block_wins reducer

def block_node(s):  return {"decision": "block"}   # fraud gate fires
def reason_node(s): return {"decision": "allow"}   # overwrites silently
def pay_node(s):    assert s["decision"] == "allow"  # always passes
g.add_edge("block", "reason"); g.add_edge("reason", "pay")
\end{lstlisting}

\noindent Declaring \texttt{decision: Annotated[str, block\_wins]} makes the same graph safe
without touching node code: the reducer, not the developer, enforces the invariant.
LangGraph's \texttt{interrupt\_before} does not close this: it presents graph state to a
reviewer, not wallet state or idempotency results, so it carries no protocol-validity
signal and composes with rather than substitutes for a fraud-semantic gate.

\begin{definition}[Structural Atomic Gate and CAR]
\label{def:gate}
A node $G$ is a \emph{structural atomic gate} iff (a)~$G$ executes synchronously before any
consequential action reaches the rail; (b)~$G$ emits a signed \emph{Commerce Action Receipt}
(CAR) $(\textit{id},\textit{action},\textit{verifications},\textit{decision},\textit{level},\textit{anchoring},\textit{sig},\textit{ts})$
carrying a \texttt{final\_decision}, where \textit{verifications} lists
$(\textit{verifier\_id},\textit{score},\textit{confidence},\textit{uncertainty})$ and
\textit{anchoring} $=(\textit{log\_hash},\textit{prev\_hash},\textit{merkle\_root},\textit{chain\_tx})$;
(c)~no graph path reaches the rail bypassing \texttt{final\_decision}; and (d)~the CAR is
persisted to the CARLog before the rail executes. These formalize the 2PC coordinator role:
no participant commits without coordinator approval, durably logged first.
\end{definition}

The ACP ceremony node satisfies Definition~\ref{def:gate}: routing is synchronous
($<$10\,ms), the CAR is signed and logged before payment routing, and routing is conditional
on \texttt{final\_decision}. With Theorem~\ref{thm:type_safety}, this yields Effectiveness:
a block from any tier cannot be overwritten downstream.

\subsection{Atomicity: Risk-Adjusted Commitment Protocols}
\label{sec:tiers}

\paragraph{Atomicity via ceremony tiers.} Commitment protocol is selected by risk.
$L_0$ (\texttt{FIND}, \texttt{QUOTE}) produces no irreversible state and needs none.
$L_1$ (\texttt{COMMIT}, amt~$<$~\$50, unc~$<$~0.25) uses Saga: each action carries a
compensating \texttt{VOID}. $L_2$ (\$50--\$5k, or unc~$\ge$~0.25) uses 2PC-lite: the CAR
coordinator requires prepare acknowledgment before commit, with Merkle-batch anchoring.
$L_3$ (amt~$>$~\$5k, unc~$\ge$~0.45, or verifier disagreement) specifies 2PC+consensus with
on-chain anchoring and is specification-only (\S\ref{sec:discussion_limits}). A
RESERVE-then-failure at any tier triggers a compensating VOID, giving Atomicity by construction.

\paragraph{Traceability via the CARLog.} Adapting Certificate Transparency~\cite{ct}, entry
$n$ stores $(\textit{seq}, \text{SHA256}(\textit{CAR}), \textit{prev\_hash}=\text{SHA256}(\textit{entry}_{n-1}), \textit{entry\_sig})$.
Modifying $\textit{entry}_k$ changes its hash, hence $\textit{entry}_{k+1}.\textit{prev\_hash}$,
breaking every subsequent signature; a chain scan detects tampering in $O(N)$. Batching every
$B$ CARs into a Merkle tree and publishing the root yields $O(\log B)$ inclusion proofs that a
third party can verify without trusting the operator.

\paragraph{Verification depth via the confidence cascade.}
\label{sec:cascade}
The cascade implements selective prediction~\cite{selective}: Tier~1 (LocalRules, confidence
$1.0$) blocks on a rule hit and otherwise delegates; Tier~2 (BehavioralML) attaches a
boundary-proximity uncertainty ($0.30$ if $|\text{score}-\text{threshold}|<0.12$, else $0.08$)
and applies the veto of Proposition~\ref{prop:veto}; Tier~3 (LLM verifier) handles the
$<$3\% of actions with high uncertainty. By Theorem~\ref{thm:type_safety} the cascade is
monotone: a later tier can add a flag or block but never retract one. These uncertainty
estimates are heuristic, not conformal intervals~\cite{conformal}
(\S\ref{sec:discussion_limits}). Appendix~\ref{app:impl} gives tier thresholds and CARLog
parameters in full.

\section{Experiments}
\label{sec:experiments}

We measure all five transaction properties across 900 synthetic commerce sessions (seeds 42--46, 100 sessions per attack type plus 100 clean sessions per seed). The companion benchmark paper~\cite{acpbench} describes the full data generation protocol.

\paragraph{Structural coverage and formal properties.}
CAO catches 300/300 definitional $\mathcal{M}$-violation events (B1, B7, B8-clear), gates
100\% of 2010 COMMIT/RESERVE actions, and logs every action---expected from the design and
reported as a correctness baseline, not a claim, since the $\mathcal{M}$-check is a DFA
simulation and is definitionally correct on grammar violations. Behavioral attacks (B2--B6)
satisfy every transition and field constraint, so it cannot see them; they are the cascade's
problem. On the formal properties: across 1,000 random permutations of verifier order,
session-level decisions are identical in 100\% of cases (Theorem~\ref{thm:type_safety},
S3); \texttt{log.verify()} detects 10/10 byte-flip corruptions in a 172-entry CARLog; and
Merkle inclusion proofs verify for 172/172 entries.

\subsection{Fraud Detection as Measured Consequence}

A system that enforces Execution Correctness and Verification is also structurally resistant to protocol-layer attacks. We verify this on ACP-Bench~\cite{acpbench} and report the results as a secondary finding. The fraud detection improvement is not a separate design goal---it is a measured consequence of satisfying the EAVT properties.

\begin{table}[h]
\caption{ACP-Bench results. ML-only and ACP rows are mean$\pm$std over seeds 42--46 on the core
8-attack benchmark (520 sessions), seed=42 in parentheses; the LLM-judge row is seed=42 over all
10 attacks. \emph{Behavioral recall} (B2--B6, B7s, B8-marginal/B8s) is the research metric;
structural recall (B1/B7/B8-clear) follows from the DFA check and is not the contribution.
Total FPR (block+flag on clean) is identical for ACP and ML-only at every seed.}
\label{tab:results}
\centering
\footnotesize
\begin{tabular}{lccccc}
\toprule
Approach & F1 & Behavioral Recall & Hard FPR & EscRate & \$/1k \\
\midrule
No verification / schema only & 0.000 & 0.000 & 0.000 & 0.000 & $<$\$0.01 \\
Rules only & $0.573{\pm}0.008$ & $0.000$ & $\mathbf{0.000}$ & $\mathbf{0.000}$ & $<$\$0.01 \\
LLM judge, Sonnet (L0 ceil.) & n/a & 0.610 & 0.000 & 0.000 & \$3.00 \\
ML only & $0.712{\pm}0.012$ & $0.740{\pm}0.021$ & $0.090{\pm}0.028$ & $0.504{\pm}0.034$ & $<$\$0.01 \\
\textbf{ACP (CAO)} & $\mathbf{0.792{\pm}0.006}$ & $\mathbf{0.880{\pm}0.009}$ & $\mathbf{0.000{\pm}0.000}$ & $0.594{\pm}0.028$ & $<$\$0.01 \\
\bottomrule
\end{tabular}
\end{table}

\textbf{The research result is not the rules column.} Rules-only reaches 0\% behavioral recall
at 0\% hard FPR---expected from a rule checker, not a finding. B7s and B8s confirm this: rules
catch 0\% of both by construction (allowlisted MCC; individual amounts within ceilings), yet
ML-only reaches 82.5\% and 100\%. The L0 ceiling---a full-session Claude Sonnet~4.6 judge,
measured rather than approximated---achieves 61\% recall at 0\% hard FPR, strong on
text-visible semantics (B2/B6/B8s at 100\%, B7s at 95\%) but limited by absent L1 state
(B4/B5/B8 at 10--25\%). \textbf{The finding} is the cascade composition: ML-only gives 74\%
behavioral recall at 9\% hard FPR, while ACP gives 88\% at 0\%, surpassing the L0 ceiling by
27 points at 300$\times$ lower cost. The gain (recall $+14$ points, hard FPR $-9$ points) comes
from two mechanisms: rules cover structural cases ML misses through behavioral feature
blindness, and the rules-certainty veto (Proposition~\ref{prop:veto}) converts ML's uncertain
blocks on legitimate high-value sessions into auditable flags. Total FPR is unchanged, so the
operational benefit is the elimination of auto-blocks, not a smaller review queue---consistent
with Corollary~\ref{cor:precision}, which locates escalation in the verifiers rather than the reducer.

\section{Limitations}
\label{sec:discussion_limits}

\textbf{Heuristic uncertainty, not conformal prediction.} The BehavioralML uncertainty
estimates are boundary-proximity heuristics (Appendix~\ref{app:impl}), not conformal
intervals~\cite{conformal}: no marginal coverage guarantee on a held-out calibration set,
no account of covariate shift between training (seed=7, $N{=}500$) and deployment, and
thresholds tuned empirically rather than derived from a coverage level. The
selective-prediction guarantee~\cite{selective} that \S\ref{sec:cascade} appeals to
therefore holds only informally. Note the limitation is confined to the \emph{verifiers}:
by Corollary~\ref{cor:precision} the reducer contributes no spurious escalation.

\textbf{Escalation rate.} At 0\% hard FPR all clean-session false positives are soft flags
(EscRate 56\%): high-value travel amounts score above the ML threshold while LocalRules
score$=0$, so the veto downgrades block$\to$flag. Reducing this without sacrificing recall
needs per-persona thresholds.

\textbf{Incomplete coverage.} B9 (revert-grant) needs post-settlement rail state and is
deferred; B4 (68\%) needs an agent-ID registry; B5 (70\%) uses ML anomaly where a settled-key
registry is correct; B7s (82.5\%) needs merchant-domain knowledge. The $L_3$ consensus tier is
specified but not implemented.

\section{Conclusion}

Whether an orchestration framework can hold a safety decision is decided by the algebra of its
state-channel reducers, and nothing else. Theorem~\ref{thm:general_char} characterizes the safe
reducers exactly---the semilattice operations dominating the lattice join---so the join is the
unique pointwise-minimal one, and last-writer-wins is unsafe on every non-trivial lattice.
The practical consequence is sharp: LangGraph's default semantics admit a bypass of \emph{any}
fraud gate placed in the graph (Theorem~\ref{thm:lww_unsound}), a structural impossibility
rather than a configuration mistake, and Theorem~\ref{thm:delegation} gives the one
condition---top-preserving embeddings---under which delegating agent populations stay
fail-closed. Instantiated as CAO, this yields all five EAVT properties; the cascade
operationalizing Verification reaches 88\% behavioral recall at 0\% hard FPR and $<$\$0.01/1k
against a 61\% L0 ceiling at \$3/1k, with B4/B5 open at 68--70\% pending cross-session
infrastructure. As agents take on consequential authority, the question these results
answer---what must be true of a framework before a safety decision survives execution---stops
being optional.

\begin{thebibliography}{99}
\small

\bibitem{saga} Garcia-Molina, H., \& Salem, K. (1987). Sagas. \emph{ACM SIGMOD Record}, 16(3), 249--259.

\bibitem{2pc} Gray, J.\ N. (1978). Notes on data base operating systems. \emph{Operating Systems: An Advanced Course}, 393--481.

\bibitem{iso8583} ISO/IEC 8583:2021. \emph{Financial transaction card originated messages: Interchange message specifications}. ISO.

\bibitem{ct} Laurie, B., Langley, A., \& Kasper, E. (2013). Certificate Transparency. \emph{RFC~6962}. IETF.

\bibitem{react} Yao, S., et al.\ (2022). ReAct: Synergizing reasoning and acting in language models. \emph{ICLR 2023}.

\bibitem{langgraph} LangChain Inc.\ LangGraph. \url{https://langchain-ai.github.io/langgraph/}, 2024.

\bibitem{pregel} Malewicz, G., et al.\ (2010). Pregel: A system for large-scale graph processing. \emph{ACM SIGMOD}, 135--146.

\bibitem{autogen} Wu, Q., et al.\ (2023). AutoGen: Enabling next-gen LLM applications via multi-agent conversation. \emph{arXiv:2308.08155}.

\bibitem{openai_agents} OpenAI. OpenAI Agents SDK. \url{https://platform.openai.com/docs/guides/agents}, 2025.

\bibitem{x402} Li, et al.\ ``Security Analysis of the x402 AI Payment Protocol.'' \emph{arXiv:2605.11781}, 2026.

\bibitem{crdts} Shapiro, M., et al.\ (2011). Conflict-free replicated data types. \emph{SSS 2011}, 386--400.

\bibitem{typestate} Strom, R.\ E., \& Yemini, S. (1986). Typestate: A programming language concept for enhancing software reliability. \emph{IEEE Trans.\ Software Eng.}, 12(1), 157--171.

\bibitem{sessiontypes} Honda, K., Vasconcelos, V.\ T., \& Kubo, M. (1998). Language primitives and type discipline for structured communication-based programming. \emph{ESOP 1998}, 122--138.

\bibitem{selective} Geifman, Y., \& El-Yaniv, R. (2017). Selective prediction in deep neural networks. \emph{NeurIPS 2017}.

\bibitem{conformal} Angelopoulos, A.\ N., \& Bates, S. (2022). A gentle introduction to conformal prediction. \emph{arXiv:2107.07511}.

\bibitem{stripe} Stripe Inc.\ ``PaymentIntent object.'' \url{https://stripe.com/docs/api/payment_intents}, 2024.

\bibitem{acpbench} [Authors]. ACP-Bench: Measuring Safety at the Payment-Protocol Layer of Agentic Commerce. \emph{AIWILD @ NeurIPS 2026}. (companion paper)


\end{thebibliography}

\section*{Broader Impact}

This framework enables compliance tooling for agentic payment systems, reducing fraud while
improving precision. Publishing the automaton grammar formalizes what ISO~8583 has enforced for
four decades and so discloses no novel attack information; Theorem~\ref{thm:general_char}
further shows no weaker typing discipline can be substituted to circumvent the gate. The
released artifact is the verifier, not exploit code.

\begin{ack}
This work was conducted without external funding. Competing interests: none.
\end{ack}

\section*{Reproducibility Statement}

All formal results are self-contained in the text and Appendix~\ref{app:general_char}; tier and
cascade parameters are in Appendix~\ref{app:impl}. The LangGraph counterexample is 15 lines of
Python reproducible with \texttt{pip install langgraph}. Empirical results reproduce via the
companion benchmark~\cite{acpbench}:
\texttt{python -m benchmark.run\_experiment --seed \{42..46\} --no-sota}.
$L_0$--$L_2$ ceremony tiers, the hash-chained CARLog, and Merkle batch anchoring are
implemented; $L_3$ consensus is specification only.


\appendix

\section{Formal Proofs}

\subsection{Proof of Theorem~\ref{thm:general_char} (Characterization of safe reducers)}
\label{app:general_char}

Throughout, $r$ is a channel reducer on a bounded safety lattice $L$, so $r(\bot,x)=x$, and
$D_i$ denotes the fold of a trace $b_1,\dots,b_n$.

\textbf{($\Leftarrow$).} Assume $r$ is associative, commutative, idempotent, and $x\sqcup y\le r(x,y)$.
(S3): associativity and commutativity make the fold of any trace equal to the $r$-product of its
multiset, which is permutation-invariant. (S4): idempotence makes the $r$-product of a multiset
equal to that of its underlying set, so duplicates do not change $D_n$. (S2): $D_{i-1}\le
D_{i-1}\sqcup b_i\le r(D_{i-1},b_i)=D_i$ and likewise $b_i\le D_i$. (S1): if $b_i=\top$ then by
(S2) $\top\le D_i$, so $D_i=\top$; by (S2) again $D_j\ge D_i=\top$ for all $j\ge i$.

\textbf{($\Rightarrow$).} Assume (S1)--(S4). Every $x\in L$ is a reachable state (trace $(x)$ gives
$D_1=r(\bot,x)=x$), so trace conditions yield identities on all of $L$.
\emph{Commutativity}: traces $(x,y)$ and $(y,x)$ give $D_2=r(x,y)$ and $r(y,x)$; (S3) equates them.
\emph{Associativity}: traces $(x,y,z)$ and $(y,z,x)$ give $r(r(x,y),z)$ and $r(r(y,z),x)=r(x,r(y,z))$
(by commutativity); (S3) equates them.
\emph{Idempotence}: traces $(x)$ and $(x,x)$ give $x$ and $r(x,x)$; (S4) equates them.
\emph{Domination}: (S2) on trace $(x,y)$ gives $x=D_1\le D_2=r(x,y)$ and $y=b_2\le D_2=r(x,y)$; so
$r(x,y)$ is an upper bound of $\{x,y\}$ in $(L,\le)$ and therefore $x\sqcup y\le r(x,y)$.

\textbf{Minimality.} $\sqcup$ is a semilattice operation dominating itself, hence safe by
($\Leftarrow$). For any safe $r$, ($\Rightarrow$) gives $\sqcup\le r$ pointwise. If $r\ne\sqcup$
then $r(x,y)>x\sqcup y$ for some pair, which is a strict escalation over $\sqcup$ on that pair
with no gain in (S1)--(S4). \qed

\subsection{Proof of Theorem~\ref{thm:delegation} (Delegation closure)}

Write $D^P$ for the parent channel state and $d^S$ for the sub-agent's final state, written
back through $\iota:L_S\to L_P$ at step $i$.
(i)~If $D^P_{i-1}=\top_P$ then $D^P_i=r_P(\top_P,\iota(d^S))=\top_P$ by (S1) for $r_P$,
whatever the sub-agent reports.
(ii)~If $d^S=\top_S$ and $\iota(\top_S)=\top_P$, then $D^P_i=r_P(D^P_{i-1},\top_P)=\top_P$ by
(S1). Conversely suppose $\iota(\top_S)=t<\top_P$. Take the trace with $D^P_{i-1}=\bot_P$ and
$d^S=\top_S$: then $D^P_i=r_P(\bot_P,t)=t\ne\top_P$, so the sub-agent's block is lost and the
composite is not fail-closed. Hence top-preservation is necessary as well as sufficient.
Finally, if $\iota$ is monotone then $a\le_S b\Rightarrow\iota(a)\le_P\iota(b)$, so the
non-regression (S2) established inside $S$ is reflected in $P$ and severity order is preserved
across the boundary. \qed

\subsection{Proof of Proposition (CTS channels are CRDTs)}

The \texttt{block\_wins} reducer is a join-semilattice on the lattice $\texttt{allow} < \texttt{flag} < \texttt{block}$: associative (three-element total order is associative), commutative (lattice join is symmetric), and idempotent ($f(x,x) = x$ for all $x$). The \texttt{operator.add} reducer over \texttt{Decimal} is a join-semilattice on non-negative reals with standard $\leq$. The \texttt{frozenset.union} reducer is a join-semilattice on the powerset lattice with subset ordering. All three are CRDTs by the CRDT join-semilattice characterization~\cite{crdts}.

\subsection{Proof of Proposition~\ref{prop:veto} (Rules-Certainty Veto)}

Let $b_i \in D = \{\texttt{allow}, \texttt{flag}, \texttt{block}\}$ be the \emph{action-level} output of the ACP node at step $i$ (after veto adjustment if applicable). Let $D_i = \texttt{block\_wins}(D_{i-1}, b_i)$ be the session-level decision after step $i$ with $D_0 = \texttt{allow}$.

\textbf{Veto trigger}: fires at step $i$ when (1)~LocalRules score$=0$ (no deterministic violation detected for action $i$); (2)~ML decision$_{\mathrm{raw}} = \texttt{block}$; (3)~uncertainty $> 0.20$. The veto adjusts $b_i \leftarrow \texttt{flag}$ before the session-level reducer runs.

\textbf{Key observation}: condition (1) means LocalRules produced decision$=\texttt{allow}$ for action $i$ specifically. It says nothing about prior actions $j < i$. The session-level state $D_{i-1}$ may already be \texttt{block} from a prior step.

\textbf{Case 1} ($D_{i-1} = \texttt{block}$): $D_i = \texttt{block\_wins}(\texttt{block}, \texttt{flag}) = \texttt{block}$. The veto cannot reduce a prior block. (Follows from Theorem~\ref{thm:type_safety}.)

\textbf{Case 2} ($D_{i-1} < \texttt{block}$): $D_i = \texttt{block\_wins}(D_{i-1}, \texttt{flag}) = \texttt{flag}$ (since \texttt{flag} $>$ \texttt{allow} in the order and $D_{i-1} \in \{\texttt{allow}, \texttt{flag}\}$). The veto correctly converts an uncertain ML block to a reviewable flag without triggering an auto-block.

In both cases the veto does not violate Theorem~\ref{thm:type_safety}: it operates at the action level before the session-level join, and the join remains monotone. \qed

\section{Implementation Parameters}
\label{app:impl}

\textbf{Ceremony tiers.} $L_0$: \texttt{FIND}, \texttt{QUOTE}; no commitment protocol, no
anchoring. $L_1$: \texttt{COMMIT} with amt~$<$~\$50 and uncertainty~$<$~0.25; Saga with
compensating \texttt{VOID}; local hash-chain anchoring. $L_2$: \texttt{COMMIT} with amt in
\$50--\$5k or uncertainty~$\ge$~0.25; 2PC-lite with prepare acknowledgment; Merkle-batch
anchoring. $L_3$: amt~$>$~\$5k, uncertainty~$\ge$~0.45, or verifier disagreement;
2PC+consensus with on-chain anchoring---\emph{specification only}, not implemented.

\textbf{Cascade thresholds.} Tier~1 LocalRules reports confidence $1.0$, uncertainty $0.0$;
a rule hit blocks without delegation, otherwise it delegates. Tier~2 BehavioralML reports
confidence $1-\hat u$ with $\hat u = 0.30$ when $|\text{score}-\text{threshold}|<0.12$ and
$0.08$ otherwise, plus additive boosts for cold-start and OOD patterns; the rules-certainty
veto fires when LocalRules score$=0$ and $\hat u > 0.20$. Tier~3 (LLM verifier) is invoked
when $\hat u > 0.35$, which occurs on $<$3\% of actions.

\textbf{CARLog.} Entry $n$ is
$\{\textit{seq}: n,\ \textit{car\_hash}: \text{SHA256}(\textit{CAR}),\ \textit{prev\_hash}: \text{SHA256}(\textit{entry}_{n-1}),\ \textit{entry\_sig}: \text{HMAC-SHA256}(\cdot)\}$.
Merkle batches use $B{=}128$ CARs per tree.

\section{Extensions}

\textbf{A2A delegation and sub-agent scope.} Theorem~\ref{thm:delegation} gives the safety
condition for delegation; the CAR provenance model additionally tracks the full delegation
chain, with a sub-agent's \texttt{COMMIT} CARed under the delegating agent's wallet and
explicit scope constraints.

\textbf{Financial trading and crypto.} The automaton $\mathcal{M}$ and CTS framework are asset-class agnostic. Equity trading agents have an analogous lifecycle (\texttt{QUOTE\_PRICE}$\to$\texttt{PLACE\_ORDER}$\to$\texttt{FILL}$\to$\texttt{CONFIRM}$\to$\texttt{SETTLE}) with a different field-constraint vocabulary (instrument allowlist, position size, leverage ratio). Crypto agents require \texttt{Reversibility=NONE} at \texttt{COMMIT} time and a chain-ID field for cross-chain replay prevention.

\textbf{Observer mode.} An observer-mode extension produces risk scores and chargeback-evidence packages without inline interception, for deployments where intermediation is not architecturally feasible.

\section*{NeurIPS Paper Checklist}

\begin{enumerate}

\item {\bf Claims}
    \item[] Question: Do the main claims made in the abstract and introduction accurately reflect the paper's contributions and scope?
    \item[] Answer: \answerYes{}
    \item[] Justification: The abstract distinguishes two types of claims: (1) structural coverage metrics (100\% M-check catch rate, FPR~=~0\% by DFA construction; 100\% CARLog traceability, architectural guarantee) and (2) probabilistic fraud detection metrics (88.0\%$\pm$0.9\% recall, 0.0\%$\pm$0.0\% hard FPR, 59.4\% escalation rate --- multi-seed over 8-attack benchmark, matched exactly in Table~\ref{tab:results}). B7s/B8s per-attack results and the LLM-judge baseline (61\% recall) are seed=42 supplementary results, labeled as such in \S6. The distinction between structural coverage claims and probabilistic detection metrics is elaborated in \S6. The Commerce State Automaton is formally defined (Definition~3.0) with proofs provided for all three theoretical results. No L3 implementation is claimed.

\item {\bf Limitations}
    \item[] Question: Does the paper discuss the limitations of the work performed by the authors?
    \item[] Answer: \answerYes{}
    \item[] Justification: \S\ref{sec:discussion_limits} covers: (1) heuristic uncertainty estimates are boundary-proximity heuristics, not conformal-prediction intervals---the exact conditions required for marginal coverage are stated; (2) 56\% escalation rate (soft-flag FPR; hard FPR is 0.0\%) is a per-persona spend-ceiling overlap limitation, not a calibration artifact; (3) B7s (MCC spoof) at 82.5\% requires merchant-domain knowledge beyond the current allowlist; (4) B4/B5 at 68--70\% require cross-session registries; (5) L3 ceremony tier (ZK proof, selective-disclosure credential) is unimplemented; (6) the commerce automaton is validated against ISO~8583, x402, and Stripe PaymentIntent but not against other payment networks.

\item {\bf Theory assumptions and proofs}
    \item[] Question: For each theoretical result, does the paper provide the full set of assumptions and a complete proof?
    \item[] Answer: \answerYes{}
    \item[] Justification: The main results are stated with explicit assumptions and complete proofs. Theorem~\ref{thm:general_char} (characterization of safe reducers on any finite bounded lattice: safe iff semilattice operation dominating the join; join is the unique minimal safe reducer) has a proof sketch in \S\ref{sec:general_typing} and a full proof in Appendix~\ref{app:general_char}; its assumptions (finite bounded lattice, reducer with $\bot$ as unit, trace semantics S1--S4) are stated in Definitions preceding it. Theorem~\ref{thm:delegation} (delegation closure: composite is fail-closed iff the lattice embedding is top-preserving) is proved in full in \S\ref{sec:general_typing}. Theorems~\ref{thm:type_safety}, \ref{thm:lww_unsound}, \ref{thm:char} (commerce instantiation) are proved in \S\ref{sec:cts}. Proposition~\ref{prop:decidability} (decidability, $O(n)$) and the tamper-evidence and CRDT propositions are proved in the text and appendix. Where a result is a direct specialization (e.g., Theorem~\ref{thm:char} from Theorem~\ref{thm:general_char}), the paper says so explicitly.

\item {\bf Experimental result reproducibility}
    \item[] Question: Does the paper fully disclose all information needed to reproduce the main experimental results?
    \item[] Answer: \answerYes{}
    \item[] Justification: The Reproducibility Statement gives the exact command (\texttt{python -m benchmark.run\_experiment --seed 42 --no-sota}), compute environment (CPU only, $<$30 seconds/seed), and JSON output field names. Markov parameters, amount distributions, perturbation $\sigma$, training seed, threshold formula, and veto conditions are fully specified in the companion AIWILD paper (cited), which contains the benchmark details. All automaton transitions are fully enumerated in Table~1.

\item {\bf Open access to data and code}
    \item[] Question: Does the paper provide open access to data and code?
    \item[] Answer: \answerYes{}
    \item[] Justification: ACP reference implementation and benchmark harness are released as open-source (anonymized URL in supplemental material; full URL in camera-ready). The benchmark is fully deterministic given \texttt{--seed}. The formal automaton definition (Definition~3.0) and CTS channel specification (\S4) are fully self-contained and implementation-independent.

\item {\bf Experimental setting/details}
    \item[] Question: Does the paper specify all training and test details necessary to understand the results?
    \item[] Answer: \answerYes{}
    \item[] Justification: Training data: seed=7, $N=500$ sessions. Test data: seeds 42--46, $N=520$ sessions each. ML threshold: $\mu + 1.1\sigma$ of clean training scores. Veto: score$=0$ and uncertainty$>0.20$. Uncertainty: boundary-proximity heuristic (0.30 if $|$score$-$threshold$|<0.12$, else 0.08). All specified in companion AIWILD paper (\S3 and Appendix~A).

\item {\bf Experiment statistical significance}
    \item[] Question: Does the paper report error bars or appropriate information about statistical significance?
    \item[] Answer: \answerYes{}
    \item[] Justification: All multi-seed results are reported as mean $\pm$ standard deviation across 5 independent seeds (42--46). Standard deviations are population standard deviations, not standard errors---stated in the table caption. The five seeds are fully independent (separate data generation). Rates are bounded well away from 0 and 1 in the reported range, so symmetric intervals are appropriate.

\item {\bf Experiments compute resources}
    \item[] Question: Does the paper provide sufficient information on compute resources needed to reproduce the experiments?
    \item[] Answer: \answerYes{}
    \item[] Justification: CPU only; no GPU required. $<$30 seconds per seed; $<$3 minutes for all five seeds. The Tier~3 LLM Verifier invoked on $<$3\% of actions in the full pipeline is bypassed by \texttt{--no-sota} in all reported experiments. Stated in the Reproducibility Statement.

\item {\bf Code of ethics}
    \item[] Question: Does the research conform with the NeurIPS Code of Ethics?
    \item[] Answer: \answerYes{}
    \item[] Justification: This work provides formal safety guarantees and defensive tooling for financial AI systems. All benchmark data is synthetic. No production payment credentials or working exploit code are released. The attack taxonomy is presented at the conceptual level required to motivate the defensive framework.

\item {\bf Broader impacts}
    \item[] Question: Does the paper discuss both potential positive and negative societal impacts?
    \item[] Answer: \answerYes{}
    \item[] Justification: See the Broader Impact section following the Conclusion. Positive: formal safety guarantees at the protocol layer (L1) reduce the window for financial fraud in agentic systems; the decidability result enables static compliance tooling. Negative: publication of the automaton model could inform adversarial agents that probe state transition boundaries; mitigated by the verifier being released alongside the taxonomy.

\item {\bf Safeguards}
    \item[] Question: Does the paper describe safeguards for responsible release of data or models with high risk for misuse?
    \item[] Answer: \answerYes{}
    \item[] Justification: The formal framework defines safety properties; the release includes the verifier (defensive tool) and the synthetic generator, not production payment attack implementations. No live API credentials, no production merchant exploit code, and no B5 settled-key replay implementations are released.

\item {\bf Licenses for existing assets}
    \item[] Question: Are creators of assets used in the paper properly credited and license terms mentioned?
    \item[] Answer: \answerYes{}
    \item[] Justification: LangGraph~\cite{langgraph} (Apache 2.0); ISO~8583 is a published standard (cited); x402~\cite{x402} (open protocol, arXiv); Stripe~\cite{stripe} (public API documentation). All appear in the References section.

\item {\bf New assets}
    \item[] Question: Are new assets introduced in the paper well documented?
    \item[] Answer: \answerYes{}
    \item[] Justification: The Commerce State Automaton (Definition~3.0), CTS channel type (Definition~4.1), and CARLog (Definition~4.2) are formally defined with full enumeration of states, transitions, and operators. The ACP reference implementation is documented in \S4 and released open-source under MIT license.

\item {\bf Crowdsourcing and research with human subjects}
    \item[] Question: Does the paper include full details of crowdsourcing or human subjects research?
    \item[] Answer: \answerNA{}
    \item[] Justification: All experimental data is synthetically generated. No human subjects, no crowdsourcing, no real user data.

\item {\bf Institutional review board (IRB) approvals}
    \item[] Question: Does the paper describe potential risks to study participants and IRB approvals?
    \item[] Answer: \answerNA{}
    \item[] Justification: No human subjects research. Not applicable.

\item {\bf Declaration of LLM usage}
    \item[] Question: Does the paper describe the usage of LLMs if it is an important, original, or non-standard component of the core methods?
    \item[] Answer: \answerYes{}
    \item[] Justification: Two LLM uses are disclosed. (1) The Tier~3 LLM Verifier (\S4) uses \texttt{claude-haiku-4-5-20251001} for high-uncertainty actions ($<3\%$ of actions), described in \S4 (``Ceremony Tiers and Escalation Protocol''). (2) The full-session LLM-as-judge baseline (\S6, Table~\ref{tab:results}) uses \texttt{claude-sonnet-4-6} via AWS Bedrock to establish the L0 detection ceiling (61\% recall, seed=42)---an evaluation component, not a method component. LLMs are \emph{not} used in the formal automaton, the CTS channel type, or the CARLog---those are provably correct by construction. LLMs were also used for writing assistance in preparing this manuscript, which does not require declaration per the NeurIPS handbook.

\end{enumerate}


\end{document}
